% ---------------------------------------------------------------------------------------
% MODELO DE TCC - INSTITUTO FEDERAL DE SERGIPE (IFS)
% ---------------------------------------------------------------------------------------
% Autor do Modelo: Mikael dos Santos Gois (mikaelgois1@gmail.com)
% Baseado nos modelos de: 
%   - Diogo D. Moreira (TCC, Dissertação e Tese)
%   - DECOM/UFOP (Monografia de Ciência da Computação)
% Modelo Base: abnTeX2 adaptado para as normativas internas de TCC do IFS (Art. 31 e 34)
% Versão: 1.0 (2026)
% ---------------------------------------------------------------------------------------
% Este arquivo faz parte de um modelo opensource para estudantes do IFS.
% Não remova os créditos dos autores originais e do modelo base.
% Este modelo é fornecido "no estado em que se encontra", sem garantias de qualquer tipo.
% ---------------------------------------------------------------------------------------

% Apêndice:
% Elemento opcional. O texto deve ser precedido da palavra APÊNDICE, identificado por letras maiúsculas
% consecutivas, travessão e pelo respectivo título. Na identificação dos apêndices, quando esgotadas
% as letras do alfabeto, devem ser utilizadas letras maiúsculas dobradas. O destaque tipográfico (negrito,
% itálico ou sublinhado) deve ser o mesmo da seção primária. - NBR 14724:2024 (4.2.3.3)

% EXEMPLO
% APÊNDICE A — AVALIAÇÃO NUMÉRICA DE CÉLULAS INFLAMATÓRIAS

\begin{apendicesenv}
  \partapendices % Indica o início dos Apendices

  \chapter{Diferença entre Anexo e Apêndice}
  Os apêndices ``são textos ou documentos elaborados pelo autor, a fim de complementar sua argumentação, sem
  prejuízoda unidade nuclear do trabalho'' \cite{NBR14724:2024}. Podem ser incluídos nos apêndices:
  os questionários da pesquisas, as tabulação de dados, ilustrações e outros documentos que necessariamente
  foram preparados pelo autor. Já os anexos, em conformidade com a norma \cite{NBR14724:2024} são
  ``textos ou documentos não elaborados pelo autor, que servem de fundamentação, comprovação e ilustração'',
  como por exemplo leis, ilustrações, demonstrações de formulas, tabulações de dados de trabalhos referenciados,
  etc.

  \chapter{Formatação}

  O texto deve ser precedido da palavra APÊNDICE, identificado por letras maiúsculas consecutivas,
  travessão e pelo respectivo título. Na identificação dos apêndices, quando esgotadas as letras do
  alfabeto, devem ser utilizadas letras maiúsculas dobradas. O destaque tipográfico (negrito, itálico
  ou sublinhado) deve ser o mesmo da seção primária. \cite{NBR14724:2024}
\end{apendicesenv}